\section{Exact calibration with a finite covariance-training sample}
\label{sec:finite-training}

This section separates three statements that have different scopes: an exact
finite-sample pivot for fixed nested linear spaces, an exact conservative
test for the nonlinear one-rate null, and a fixed-$K$ limit for the fitted
phase geometry.  Throughout, the signal average and the covariance-training
sample are independent and Gaussian.  Write
\begin{align*}
X&=\sqrt m\,(\bar Y-M)\sim N_d(0,\Sigma),\\
W&=\nu\widehat\Sigma\sim W_d(\Sigma,\nu),\\
\nu&=K-1.
\end{align*}
where the $K$ training vectors have been centered before their sample
covariance is formed.  The implementation considered here has
$\nu\ge d$, so that $\widehat\Sigma$ is nonsingular almost surely.

\begin{theorem}[finite-training nested-space pivot]
\label{thm:finite-training-pivot}
Let $\mathcal A\subset\mathcal L$ be fixed, nonrandom linear subspaces of
$\mathbb R^d$, with
\[
 p=\dim(\mathcal A),\qquad r=\dim(\mathcal L)-p,
 \qquad k=d-p.
\]
Under $M\in\mathcal A$, define
\begin{align*}
Q_{\mathcal A}&=m\inf_{a\in\mathcal A}
 (\bar Y-a)^T\widehat\Sigma^{-1}(\bar Y-a),\\
Q_{\mathcal L}&=m\inf_{\ell\in\mathcal L}
 (\bar Y-\ell)^T\widehat\Sigma^{-1}(\bar Y-\ell).
\end{align*}
Then, for every positive definite $\Sigma$ and every $M\in\mathcal A$,
\[
 F_{\rm nest}={\nu-k+1\over r}
 {Q_{\mathcal A}-Q_{\mathcal L}\over \nu+Q_{\mathcal L}}
 \sim F_{r,\nu-k+1}.
\]
Consequently this calibration is invariant under every nonsingular common
change of coordinates applied to the observation, training vectors, and the
two spaces.  Its null law depends only on $(r,k,\nu)$.
\end{theorem}

\begin{proof}
Apply the true covariance whitening map and then an orthogonal rotation.
Because the statistic is unchanged by this common transformation, take
$\Sigma=I_d$, let the first $p$ coordinates span $\mathcal A$, and let the
next $r$ coordinates complete it to $\mathcal L$.  Minimizing over the first
$p$ mean coordinates reduces the problem to an independent
$x\sim N_k(0,I_k)$ and $W_0\sim W_k(I_k,\nu)$.  Partition their coordinates
as $r+q$, where $q=k-r$, and write
\[
 W_0=\begin{pmatrix}W_{BB}&W_{BR}\\W_{RB}&W_{RR}\end{pmatrix},
 \qquad x=\begin{pmatrix}x_B\\x_R\end{pmatrix}.
\]
If $q_0=Q_{\mathcal A}/\nu$ and
$q_1=Q_{\mathcal L}/\nu$, block inversion gives
\[
 q_1=x_R^TW_{RR}^{-1}x_R,
 \quad
 q_0-q_1=a^TH^{-1}a,
\]
where
\[
 a=x_B-W_{BR}W_{RR}^{-1}x_R,
 \qquad
 H=W_{BB}-W_{BR}W_{RR}^{-1}W_{RB}.
\]
Wishart regression decomposition gives
$H\sim W_r(I_r,\nu-q)$, independently of $W_{RR}$ and the corresponding
regression coefficient.  Conditional on $(x_R,W_{RR})$,
\[
 {a\over\sqrt{1+q_1}}\sim N_r(0,I_r),
\]
and this standardized vector is independent of $H$.  The ordinary
Hotelling quadratic-form identity therefore yields
\[
 {\nu-q-r+1\over r}{q_0-q_1\over1+q_1}
 \sim F_{r,\nu-q-r+1}.
\]
Substitution of $q=k-r$ and $q_j=Q_j/\nu$ proves the result.
\end{proof}

The denominator $\nu+Q_{\mathcal L}$ is essential.  A scalar correction to
$Q_{\mathcal A}-Q_{\mathcal L}$ can correct an inverse-Wishart expectation,
but it cannot remove the finite-training random scale or its heavy tail.
Only in the special case $\mathcal L=\mathbb R^d$ is
$Q_{\mathcal L}=0$ and the result reduces directly to the usual Hotelling
$T^2$--$F$ identity.

\begin{proposition}[finite-sample conservative nonlinear-null test]
\label{prop:finite-training-omnibus}
For a candidate rate $\theta$, let $\mathcal A_\theta$ be the linear space
generated by the channel-specific intercepts and amplitudes.  Suppose its
rank is the constant $p_0$ throughout the rate search interval, and put
$k_0=d-p_0$.  Define the profiled residual
\[
 Q_{\rm prof}=m\inf_{\theta}\inf_{a\in\mathcal A_\theta}
  (\bar Y-a)^T\widehat\Sigma^{-1}(\bar Y-a).
\]
For every finite $m$ and every one-rate null mean whose true rate lies in the
search interval, the rule
\[
 {\nu-k_0+1\over \nu k_0}Q_{\rm prof}
 >F^{-1}_{k_0,\nu-k_0+1}(1-\alpha)
\]
has rejection probability at most $\alpha$.
\end{proposition}

\begin{proof}
At the true rate $\theta_0$, minimization over its $p_0$ linear coefficients
and the same whitening-and-rotation argument give
\[
 {\nu-k_0+1\over \nu k_0}Q_{\theta_0}
 \sim F_{k_0,\nu-k_0+1}.
\]
Since profiling over the rate can only decrease the residual,
$Q_{\rm prof}\le Q_{\theta_0}$ pointwise.  The rejection event is therefore
contained in the level-$\alpha$ fixed-rate event.
\end{proof}

In the exponential design of the main text, $p_0=2C$.  Thus for
$(C,n,d)=(4,12,48)$ one has $k_0=40$ and denominator degrees of freedom
$K-40$.

\begin{theorem}[scope of fitted phase calibration at fixed $K$]
\label{thm:finite-training-plugin}
Fix $d$ and $K$ with $K-1\ge d$.  Assume an interior identifiable one-rate
null, a nonzero channel-amplitude vector, constant tangent and collision
ranks, and uniform nonsingularity of their relevant Gram matrices.  Let
$\widehat{\mathcal A}$ be the tangent at the GLS null fit and let
$\widehat{\mathcal L}$ add the $r=C$ collision directions.  Form
$\widehat Q_{\mathcal A}$ and $\widehat Q_{\mathcal L}$ by the corresponding
estimated-metric projections.  Under the null, as $m\to\infty$ with $K$
fixed,
\[
 {\nu-k+1\over r}
 {\widehat Q_{\mathcal A}-\widehat Q_{\mathcal L}
  \over \nu+\widehat Q_{\mathcal L}}
 \ \Longrightarrow\ F_{r,\nu-k+1},
 \qquad k=d-(2C+1).
\]
For $(C,n,d)=(4,12,48)$ this is $F_{4,K-39}$ after multiplication by
$(K-39)/4$.
\end{theorem}

\begin{proof}[Proof sketch]
Conditional on the almost surely nonsingular training scatter, standard
fixed-dimensional nonlinear least-squares expansion makes the fitted null
tangent and collision extension converge to their population counterparts;
the difference between the fitted projection residuals and those in the
fixed tangent experiment is $o_p(1)$.  The continuous mapping theorem and
Theorem~\ref{thm:finite-training-pivot} then give the displayed unconditional
limit.  No assertion of exactness at finite $m$ is made.
\end{proof}

Conditional on the observed $\widehat\Sigma$, there is no covariance-free
null calibration in general.  Already for $d=1$, conditioning on
$\widehat\Sigma=s$ leaves a quadratic form distributed as
$(\Sigma/s)\chi_1^2$.  The exact $F$ law arises only after averaging over the
independent Wishart training sample.  Likewise, the nonlinear fitted geometry
precludes a universal exact finite-$m$ phase law: Gaussian distance to a
curved mean manifold depends on its local curvature and nuisance location.
Splitting only the covariance-training sample does not remove either issue.
Finally, this section does not assert an exact finite-training mixture-$F$
law for the one-sided cone statistic; selection of a cone face and the random
Wishart metric require a separate proof.
